% ==========================================================================
\section{Conclusão}
\label{sec:conclusao}
% ==========================================================================

Este trabalho permitiu concluir com sucesso o projeto e a análise comparativa de LNAs para aplicações
de radiofrequência, cobrindo tanto a implementação com componentes discretos de alto desempenho como
a integração em tecnologia CMOS em três nós tecnológicos distintos.

Na primeira fase, o dimensionamento do LNA com o transístor bipolar BFP420 estabeleceu os
fundamentos metodológicos do projeto RF: definição do ponto de operação estável na região ativa a
6\,GHz, análise de estabilidade com recurso ao fator de Rollet ($K$) e ao AppCAD, e síntese de redes
de adaptação por inspeção na Carta de Smith. Tanto a abordagem por componentes discretos ($LC$) como
a modelação por linhas de transmissão com \textit{stubs} demonstraram um bom casamento de impedâncias
($S_{11}$, $S_{22} < -10~\text{dB}$), tendo a transição para o Cadence permitido co-otimizar o ganho
real (10{,}15\,dB) e a figura de ruído (2{,}26\,dB) com uma precisão que o LTspice não assegura,
dado o impacto dos parasitas de layout.

A segunda fase, centrada no \textit{balun-LNA} em tecnologia CMOS com a arquitetura de
Blaakmeer~\cite{blaakmeer2008}, aprofundou os aspetos específicos do dimensionamento integrado.
Demonstrou-se que a adaptação de entrada a $50~\Omega$ pode ser obtida de forma ativa, através do
controlo de $g_m$ no estágio CG, e que o mecanismo de \textit{admittance scaling} permite reduzir a figura de 
ruído com rendimento marginal decrescente, mantendo o balanço diferencial. Os resultados a 65\,nm confirmaram
quantitativamente este comportamento: a redução de NF de $N=1$ para $N=2$ é de 16\,\%, enquanto de
$N=2{,}75$ para $N=3{,}5$ é apenas de 3{,}6\,\%, e a contribuição de $r_g$ ao nível de 7\,\%
sinaliza o limite prático de rendimento do escalamento.

O estudo dos três nós tecnológicos revelou tendências fisicamente coerentes e pedagogicamente
relevantes. Em 350\,nm, a validade das equações de canal longo garante boa previsibilidade analítica.
Em 65\,nm, os efeitos de canal curto como saturação de velocidade, degradação da mobilidade e
aumento de $r_g$ começam a limitar o desempenho, mas o modelo ainda suporta uma análise mista
analítico-numérica. Em 45\,nm, a transcondutância entra no regime de saturação de velocidade
descrito por~\eqref{eq:gm_vsat}, o ganho intrínseco colapsa para valores da ordem de 7, e a
extração numérica pelo simulador BSIM4 torna-se o único referencial fiável. A inversão de ruído
observada entre 65\,nm e 45\,nm, onde o nó mais avançado apresenta ligeiramente mais ruído total
nas condições de polarização utilizadas, é um resultado particularmente expressivo, que demonstra
que a progressão tecnológica não é automaticamente sinónimo de melhor desempenho RF sem otimização
dedicada de polarização e geometria.

Do ponto de vista das especificações do projeto, os objetivos de $S_{11} < -10~\text{dB}$,
$S_{22} < -10~\text{dB}$ e ganho $> 10~\text{dB}$ foram atingidos em 350\,nm e 65\,nm. Em 45\,nm,
a compressão do ganho intrínseco e o aumento das perdas resistivas dificultam o cumprimento
simultâneo de todas as especificações, ilustrando os desafios reais do projeto RF em nós nanométricos
avançados.

Em termos de comparação com o design bipolar, a solução CMOS integrada distingue-se pelo balun ativo
sem indutores, pela largura de banda intrínseca superior e pela compatibilidade com processos digitais
de larga escala, ao custo de maior complexidade de dimensionamento e maior sensibilidade à variabilidade
de processo nos nós mais avançados. Os dois paradigmas são complementares: o BFP420 é a escolha para
desempenho máximo em banda estreita com componentes discretos, enquanto as topologias CMOS são
adequadas para integração wideband em transceivers multistandard modernos.